\documentclass[11pt]{article}
\usepackage[utf8]{inputenc}
\usepackage{amsmath, amssymb, amsthm}
\usepackage[margin=1in]{geometry}

\theoremstyle{definition}
\newtheorem{definition}{Definition}
\newtheorem{theorem}{Theorem}
\newtheorem{lemma}{Lemma}
\newtheorem{proposition}{Proposition}

\theoremstyle{remark}
\newtheorem{remark}{Remark}

\DeclareMathOperator{\tr}{tr}
\newcommand{\dual}[2]{{}_{V'}\langle #1, #2 \rangle_{V}}

\title{Lecture 4: Stochastic Integrals and SPDEs on Hilbert Spaces}
\date{06.11.2025}
\author{Tien Anh Vu}

\begin{document}
\maketitle

This lecture completes the construction of the infinite-dimensional stochastic integral, develops its main structural properties, and then uses this machinery to formulate SPDEs on Hilbert spaces, introduce the different notions of solution, and prove existence and uniqueness of mild solutions by a fixed point argument.

\section*{1. Extending the Stochastic Integral}
We first turn the Wiener-It\^o isometry for elementary processes into an extension principle for the stochastic integral.

\subsection*{1.1. The Wiener-It\^o Isometry (Recap)}
The extension begins with the identity that makes it possible:
\[
\|I(\Phi)\|_{\mathcal M_T^2}^2=\|\Phi\|_{\mathcal E}^2.
\]
This states that the norm of the integral mapping \(I(\Phi)\) inside the space of \(L^2\)-bounded martingales \(\mathcal M_T^2\) is equal to the norm of the elementary process \(\Phi\) in its own space \(\mathcal E\). Because this mapping preserves the norm, it is an isometry, and this is exactly the property that allows us to extend the integral beyond elementary processes.

\subsection*{1.2. Step 3: Extending to the Closure}
Since the elementary space is too small for applications, the next step is to identify its closure and extend the integral continuously to that larger space.

To define this closure, we first introduce the predictable \(\sigma\)-algebra, denoted by \(\mathcal P_T\). It is defined by
\[
\mathcal P_T
:=
\sigma\bigl(\{]s,t]\times F_s:0\le s<t\le T,\ F_s\in\mathcal F_s\}\cup\{\{0\}\times F_0:F_0\in\mathcal F_0\}\bigr).
\]
This is the \(\sigma\)-algebra generated by predictable rectangles over time and probability space.
Intuitively, it collects exactly those events that can be described using information available strictly up to the current time.

An equivalent and more operational description is that \(\mathcal P_T\) is generated by all left-continuous \((\mathcal F_t)\)-adapted processes \((H_s)_{s\in[0,T]}\). Left-continuity ensures that the process depends only on strictly past information and does not look into the future.

With the predictable \(\sigma\)-algebra in place, the closure is described by
\[
\overline{\mathcal E}
=
L^2(\Omega_T,\mathcal P_T,\mathbb P_T;L_2^0)/\sim.
\]
Here
\[
\Omega_T=\Omega\times[0,T],
\qquad
\mathbb P_T=\mathbb P\otimes dt,
\]
so \(\Omega_T\) is the space of all pairs \((\omega,t)\), and \(\mathbb P_T\) is the product measure combining probability on \(\Omega\) with Lebesgue measure in time.
and
\[
L_2^0
:=
L_2(\sqrt Q\,U,H)
\]
denotes the space of all continuous linear operators mapping from \(U\) to \(H\) such that their composition with \(\sqrt Q\) is Hilbert-Schmidt. The quotient \(/\sim\) means that we identify processes that are equal almost everywhere with respect to the norm.

Because the Wiener-It\^o isometry is already known for elementary processes, the integral operator \(I\) extends continuously to this entire closure.

\subsection*{1.3. Step 4: Localization}
Even this \(L^2\)-closure remains too restrictive for many applications, so we next pass to the standard localization argument.

The original \(L^2\)-theory was built under the expectation-level integrability condition
\[
\mathbb E\int_0^T \|\Phi(s)\circ \sqrt Q\|_{L_2(U,H)}^2\,ds<\infty,
\]
but many useful integrands do not satisfy this globally. In applications, one often has only the weaker pathwise condition
\[
\int_0^T \|\Phi(s,\omega)\circ \sqrt Q\|_{L_2(U,H)}^2\,ds<\infty
\qquad \text{for almost every } \omega.
\]
Equivalently,
\[
\mathbb P\left(\int_0^T \|\Phi(s)\circ \sqrt Q\|_{L_2(U,H)}^2\,ds<\infty\right)=1.
\]
Localization allows us to define the integral for all processes \(\Phi:\Omega_T\to L_2^0\) that are adapted, \(\mathcal P_T\)-measurable, and satisfy this almost-sure integrability condition.

To return to the old \(L^2\)-framework, we localize by stopping times. Choose a sequence of stopping times \((\tau_n)\) that cuts the process off before the pathwise integral becomes too large, for example when it reaches the level \(n\). Then on each stopped interval \([0,\tau_n]\) one has
\[
\int_0^{T\wedge \tau_n} \|\Phi(s)\circ \sqrt Q\|_{L_2(U,H)}^2\,ds\le n
\qquad \text{a.s.}
\]
Taking expectation yields
\[
\mathbb E\int_0^{T\wedge \tau_n} \|\Phi(s)\circ \sqrt Q\|_{L_2(U,H)}^2\,ds\le n<\infty.
\]
Thus, after stopping, the integrand falls back into the previously constructed \(L^2\)-theory, and the stochastic integral can be defined on each localized interval.

\subsection*{1.4. Core Properties of the Stochastic Integral}
With the localized construction in place, we record two basic operational properties of the stochastic integral.

\textbf{(i) Linearity.} Suppose \(L\in L(H,\widetilde H)\) is a continuous linear operator from \(H\) into another Hilbert space \(\widetilde H\). Then
\[
L\left(\int_0^t \Phi(s)\,dW(s)\right)
=
\int_0^t L\circ\Phi(s)\,dW(s).
\]
Thus a continuous linear operator may be pulled inside the stochastic integral.

\textbf{(ii) Inner products.} Let \(f:\Omega_T\to H\) be a continuous, \((\mathcal F_t)\)-adapted process. Then the integral respects the inner product structure of \(H\). We define
\[
\int_0^t \langle f(s),\Phi(s)\,dW_s\rangle_H
:=
\int_0^t \Phi_f^*(s)\,dW_s.
\]
Here \(\Phi_f^*(s)\in L(U,\mathbb R)\) is the scalar-valued operator defined by
\[
\Phi_f^*(s)(u)=\langle f(s),\Phi(s)u\rangle_H,
\qquad u\in U.
\]
Thus the general takeaway is that the stochastic integral of the inner-product process can be rewritten as an ordinary real-valued stochastic integral with the modified integrand \(\Phi_f^*(s)\).

\section*{2. Advanced Properties of the Stochastic Integral}
Having completed the construction of the integral, we now turn to the estimates and structural properties needed later in SPDE theory.

\subsection*{2.1. The Conclusion of the Inner Product Property}
We first complete the inner-product identity from the previous section and define the associated scalar-valued operator by
\[
\Phi_f^*(s)(u)=\langle f(s),\Phi(s)u\rangle_H,
\qquad u\in U.
\]
This formally completes the transition of the deterministic process \(f(s)\) into the operator structure.

\subsection*{2.2. Property (iii): The Burkholder-Davis-Gundy (BDG) Inequalities}
We next introduce the main maximal estimate for the stochastic integral, namely the Burkholder-Davis-Gundy inequalities.

This inequality is a substantial generalization of both Doob's maximal inequality and the Wiener-It\^o isometry. For any exponent \(p\), the notes state that
\[
\mathbb E\left[\sup_{t\in[0,T]}\left\|\int_0^t \Phi(s)\,dW(s)\right\|_H^{2p}\right]
\le
C_p\,\mathbb E\left[\left(\int_0^T \|\Phi(s)\|_{L_2^0}^2\,ds\right)^p\right].
\]
On the left-hand side, we control the expected maximum of the path norm, that is, the supremum over \(t\in[0,T]\) of the \(H\)-norm of the stochastic integral, raised to the power \(2p\). The inequality says that this supremum is bounded by a constant \(C_p\) times the \(p\)-th moment of the time integral of the squared \(L_2^0\)-norm of the integrand. In this sense, the moment of the path supremum is controlled by the variation of the martingale, and extreme pathwise behavior is reduced to an ordinary Lebesgue integral.

\subsection*{2.3. Property (iv): Quadratic Variation}
After the maximal estimate, we identify the quadratic variation of the resulting martingale.

Let \(M_t:=I(\Phi)=\int_0^t \Phi(s)\,dW(s)\). Then its quadratic variation is defined by
\[
\langle M\rangle_t=\int_0^t \|\Phi(s)\|_{L_2^0}^2\,ds,
\qquad t\in[0,T].
\]
The notes explicitly remind us that the notation \(L_2^0\) refers to the Hilbert-Schmidt norm after composition with \(\sqrt Q\).

This process is characterized as the unique continuous, increasing, and \((\mathcal F_t)_{t\in[0,T]}\)-adapted process starting at zero such that
\[
\|M_t\|_H^2-\langle M\rangle_t
\]
is itself an \((\mathcal F_t)_{t\in[0,T]}\)-martingale. In this way, the quadratic variation records exactly how the variance of the stochastic integral accumulates over time.

\subsection*{2.4. Property (v): Regularity of the Stochastic Integral}
We then ask how regular the stochastic integral paths are in time.

For a parameter \(\alpha<1/2\) and an integrand \(\Phi\) belonging to the standard \(L^2\) space, the process \(M_t:=I(\Phi)\) belongs to the fractional Sobolev space \(W^{\alpha,2}([0,T],H)\), defined by
\[
W^{\alpha,2}([0,T],H)
:=
\left\{
M\in L^2([0,T],H):
\int_0^T\int_0^T
\frac{\mathbb E\|M_s-M_t\|_H^2}{|s-t|^{1+2\alpha}}
\,ds\,dt<\infty
\right\}.
\]
Thus the path regularity is measured through the finiteness of a double integral that compares the process at all time pairs \(s\) and \(t\). Since Brownian paths are nowhere differentiable, the correct regularity is described by these fractional spaces rather than by classical derivative spaces.

\subsection*{2.5. Proof Sketch: Verifying the Regularity}
To justify this regularity statement, we estimate the defining double integral directly. Let us denote
\[
X_T
:=
\int_0^T\int_0^T
\frac{\mathbb E\|M_s-M_t\|_H^2}{|s-t|^{1+2\alpha}}
\,ds\,dt.
\]
We first view the difference \(M_s-M_t\) itself as a stochastic integral over the interval between \(s\) and \(t\). Applying the Wiener-It\^o isometry to that increment then yields
\begin{align*}
\mathbb E\|M_s-M_t\|_H^2
&=
\mathbb E\left\|\int_{s\wedge t}^{s\vee t}\Phi(r)\,dW(r)\right\|_H^2 \\
&=
\mathbb E\left[\int_{s\wedge t}^{s\vee t}\|\Phi(r)\|_{L_2^0}^2\,dr\right].
\end{align*}
\[
\mathbb E[X_T]
=
\int_0^T\int_0^T
\frac{\mathbb E\left[\int_{s\wedge t}^{s\vee t}\|\Phi(r)\|_{L_2^0}^2\,dr\right]}
{|s-t|^{1+2\alpha}}
\,ds\,dt.
\]
Using the symmetry of the integration domain, this becomes
\[
\mathbb E[X_T]
=
2\int_0^T\int_0^t
\frac{\mathbb E\left[\int_s^t \|\Phi(r)\|_{L_2^0}^2\,dr\right]}
{(t-s)^{1+2\alpha}}
\,ds\,dt.
\]
By Fubini's theorem, the order of integration may be changed:
\[
\mathbb E[X_T]
=
2\int_0^T
\mathbb E\bigl(\|\Phi(r)\|_{L_2^0}^2\bigr)
\left(
\int_r^T\int_0^r \frac{1}{(t-s)^{1+2\alpha}}\,ds\,dt
\right)dr.
\]
The inner time integral depends only on \(\alpha\) and \(T\), and it is finite for \(\alpha<1/2\). Hence there exists a constant \(C(\alpha)\) such that
\[
\int_r^T\int_0^r \frac{1}{(t-s)^{1+2\alpha}}\,ds\,dt
\le C(\alpha)
\qquad \text{for all } r\in[0,T].
\]
Therefore
\[
\mathbb E[X_T]
\le
C(\alpha)\int_0^T \mathbb E\bigl(\|\Phi(r)\|_{L_2^0}^2\bigr)\,dr.
\]
Because \(\Phi\) belongs to the standard \(L^2\) space, the right-hand side is finite. Therefore the stochastic integral belongs to \(W^{\alpha,2}([0,T],H)\).

\section*{3. The General Framework for SPDEs}
With the stochastic integral in place, we now pass to its primary application, namely stochastic partial differential equations on Hilbert spaces.

\subsection*{3.1. The Cylindrical Wiener Process and Integrability}
We begin by clarifying the distinction between the \(Q\)-Wiener process and the cylindrical Wiener process, since this determines the type of noise that may drive an SPDE.

We recall the standard \(Q\)-Wiener process
\[
W_t=\sum_{n\in\mathbb N}\sqrt{\lambda_n}\,\beta_t^{(n)}u_n.
\]
As long as the covariance operator \(Q\) has finite trace, that is,
\[
\sum_{n\in\mathbb N}\lambda_n<\infty,
\]
the series defining \(W_t\) converges in \(L^2\). More precisely, for each fixed \(t\), the partial sums
\[
W_t^N:=\sum_{n=1}^N \sqrt{\lambda_n}\,\beta_t^{(n)}u_n
\]
converge in \(L^2(\Omega;U)\) to a \(U\)-valued random variable \(W_t\).

Against this background, the notes introduce the cylindrical Wiener process,
\[
W_t=\sum_{n\in\mathbb N}\beta_t^{(n)}u_n.
\]
Here the eigenvalues \(\lambda_n\) are absent, so all infinite directions carry the same variance. Consequently, this series does not converge in the original Hilbert space \(U\); it represents the idealized infinite-dimensional noise that corresponds to space-time white noise.

To integrate with respect to this noise, we recall the fundamental integrability condition:
\[
\int_0^T \|\Phi_s\circ \sqrt Q\|_{L_2(U,H)}^2\,ds<\infty.
\]
In other words, the time-integrated squared Hilbert-Schmidt norm of the integrand must be finite.
This is the condition that keeps the stochastic integral mathematically well behaved.

\subsection*{3.2. The General Framework for SPDEs}
With our noise established, we formalize the core environment for our SPDEs. Let \(U\) be the Hilbert space of the driving Wiener process, let \(H\) be the Hilbert space in which the unknown state evolves, and let \((W_t)_{t\in[0,T]}\) be a \(Q\)-Wiener process. We then write the SPDE in its abstract stochastic differential equation form on \(H\):
\[
dX_t=[AX_t+B(X_t)]\,dt+C(X_t)\,dW_t,
\]
Alongside this, we are given an initial condition
\[
X_0=\xi.
\]
The evolution of the state \(X_t\) is driven by three distinct components:
\[
AX_t,
\qquad
B(X_t),
\qquad
C(X_t)\,dW_t.
\]
The first term is the linear drift, typically an unbounded differential operator such as a Laplacian. The second term is the nonlinear drift component. The third term is the diffusion component, representing how the system reacts to the stochastic noise.

\subsection*{3.3. The Four Fundamental Assumptions (A.1--A.4)}
To ensure that this SPDE is solvable, we cannot allow the operators \(A\), \(B\), and \(C\) to behave arbitrarily. The notes therefore impose four standard assumptions.

\textbf{(A.1) The linear operator.} The pair \((A,D(A))\) generates a \(C_0\)-semigroup \((T_t)_{t\ge 0}\) on \(H\). Because \(A\) is often a spatial derivative, it is an unbounded operator and cannot be applied to every element of \(H\). Instead, it is defined on a specific domain \(D(A)\), and the semigroup it generates is used to evolve the system in time.

\textbf{(A.2) The nonlinear drift.} The mapping
\[
B:H\to H
\]
is \(\mathcal B(H)\)-measurable. This is a standard topological requirement ensuring that \(B\) is well behaved within the Borel \(\sigma\)-algebra.

\textbf{(A.3) The diffusion operator.} The mapping
\[
C:H\to L_2(U_0,H)
\]
is strongly continuous. Note that \(C\) maps the state into the space of Hilbert-Schmidt operators. Strong continuity means that for every fixed \(u\in U_0\), the map
\[
x\mapsto C(x)u
\]
from \(H\) to \(H\) is continuous.
This is the natural target space for the diffusion term because the stochastic integral is defined through the Hilbert-Schmidt condition on \(C(x)\circ \sqrt Q\).

\textbf{(A.4) The initial condition.} The initial state \(\xi\) is an \(H\)-valued random variable that is \(\mathcal F_0\)-measurable. This ensures that our starting point uses only the information available at time \(t=0\).

\subsection*{3.4. Defining the \(C_0\)-Semigroup}
Since assumption \textnormal{(A.1)} is central to the whole SPDE setup, we now recall the precise definition of a \(C_0\)-semigroup.

A family of bounded linear operators \((T_t)_{t\ge 0}\) on \(H\) is officially called a \(C_0\)-semigroup if it strictly satisfies three properties:
\begin{itemize}
\item \textbf{The identity.} At time zero, the operator does nothing:
\[
T_0=\operatorname{Id}.
\]
\item \textbf{The semigroup property.} Evolving the system for time \(s\), and then subsequently for time \(t\), is identical to evolving it directly for time \(t+s\). Mathematically,
\[
T_t\circ T_s=T_{t+s}.
\]
\item \textbf{Strong continuity.} For any fixed state \(u\in H\), the mapping
\[
t\mapsto T_tu
\]
is continuous.
\end{itemize}

Finally, we link this semigroup back to the linear operator \(A\). One says that \(A\) generates the \(C_0\)-semigroup \((T_t)_{t\ge 0}\), and the generator is defined by
\[
Ah=\lim_{t\to 0}\frac{T_th-T_0h}{t}
=
\lim_{t\to 0}\frac{T_th-h}{t}.
\]
This definition specifies both how the operator \(A\) acts and how its domain \(D(A)\) is defined, namely as the set of all \(h\in H\) for which this limit exists. In this way, the semigroup framework lets us work rigorously with unbounded operators such as the Laplacian and prepares the way for the notions of strong, mild, and weak solutions that come next.

\section*{4. Notions of Solution for SPDEs}
With the semigroup framework in place, we can now address the next question: what does it mean to solve an SPDE?

\subsection*{4.1. The Mild Solution}
We begin with the most practical and widely used concept in SPDE theory, namely the mild solution.

An \(H\)-valued predictable process \((X_t)_{t\in[0,T]}\) is called a mild solution if it satisfies the following integral equation \(\mathbb P\)-almost surely:
\[
X_t
=
T_t\xi+\int_0^t T_{t-s}B(X_s)\,ds+\int_0^t T_{t-s}C(X_s)\,dW_s.
\]
Let us dissect the mechanics of this equation. This is the infinite-dimensional analogue of the variation of constants formula. Instead of differentiating, we integrate, utilizing the \(C_0\)-semigroup \((T_t)\) as our temporal propagator.

The first term, \(T_t\xi\), propagates the initial condition forward in time. The second term convolves the nonlinear drift \(B(X_s)\) with the semigroup \(T_{t-s}\). The final term convolves the stochastic diffusion \(C(X_s)\) with the semigroup \(T_{t-s}\) with respect to the Wiener process.

A critical analytical warning appears directly below the formula. We cannot simply factor out \(T_t\) from the integrals to write an expression involving \(T_{-s}\). Standard semigroups evolve only forward in time, and the inverse operator \(T_{-s}\) is generally not well defined. Therefore the convolution \(T_{t-s}\) must remain inside the integrals, and all of these integrals must be independently well defined for the mild solution to exist.

\subsection*{4.2. The (Analytically) Strong Solution}
For comparison, we next examine the strong solution, which is the most rigid formulation.

An analytically strong solution is a predictable process \((X_t)_{t\ge 0}\) that satisfies the SPDE in its pure differential form, integrated directly:
\[
X_t
=
\xi+\int_0^t [AX_s+B(X_s)]\,ds+\int_0^t C(X_s)\,dW_s.
\]
Notice what is missing here: there is no semigroup \(T_{t-s}\) smoothing out the equation. Because the unbounded linear operator \(A\) is applied directly to the state \(X_s\), this formulation demands a severe restriction: the process \((X_t)\) must be a \(D(A)\)-valued process. In infinite dimensions this requirement is usually far too restrictive. The spatial roughness of the noise \(dW_s\) often pushes the state \(X_t\) immediately out of the domain \(D(A)\), so analytically strong solutions rarely exist for typical SPDEs.

\subsection*{4.3. The (Analytically) Weak Solution}
To bypass the strict domain restriction of the strong solution, we now turn to the weak solution.

Here we require an \(H\)-valued predictable process \((X_t)_{t\in[0,T]}\) that satisfies the equation tested against a smooth function \(\varphi\). For every \(\varphi\in D(A^*)\), we write
\[
\langle X_t,\varphi\rangle_H
=
\langle \xi,\varphi\rangle_H
+\int_0^t \langle X_s,A^*\varphi\rangle_H\,ds
+\int_0^t \langle B(X_s),\varphi\rangle_H\,ds
+\int_0^t \langle \varphi,C(X_s)\,dW_s\rangle_H.
\]
The mathematical elegance of this formulation lies in the second term on the right. We have shifted the burden of the unbounded operator \(A\) off of the rough stochastic process \(X_s\) and onto the smooth test function \(\varphi\), using the identity
\[
\langle AX_s,\varphi\rangle_H=\langle X_s,A^*\varphi\rangle_H.
\]
As long as we choose our test functions \(\varphi\) in \(D(A^*)\), the equation is well defined without forcing \(X_t\) to have the high spatial regularity required by the strong formulation.

\subsection*{4.4. Existence and Uniqueness of Mild Solutions}
Having established these three notions of solution, we focus entirely on the mild solution. To guarantee that a unique mild solution exists, we impose two strict functional bounds on the nonlinear drift \(B\) and the diffusion \(C\), denoted by \textnormal{(H.1)} and \textnormal{(H.2)}.

\textbf{(H.1) Global Lipschitz condition.}
\[
\|B(t,x)-B(t,y)\|_H+\|(C(t,x)-C(t,y))\circ \sqrt Q\|_{L_2(U,H)}
\le
L\|x-y\|_H
\]
for all \(x,y\in H\). This condition ensures that the operators do not react wildly to small changes in the state. For the diffusion term \(C\), the distance is measured by
\[
\|(C(t,x)-C(t,y))\circ \sqrt Q\|_{L_2(U,H)},
\]
which is exactly the norm that appears in the definition of the stochastic integral. This Lipschitz bound is the key condition for uniqueness.

\textbf{(H.2) Linear growth condition.}
\[
\|B(t,x)\|_H^2+\|C(t,x)\circ \sqrt Q\|_{L_2(U,H)}^2
\le
M(1+\|x\|_H^2)
\]
for all \(x\in H\). This inequality bounds the growth of the drift and diffusion and prevents the stochastic process from exploding in finite time. The linear growth condition guarantees global existence on the whole interval \([0,T]\).

\subsection*{4.5. Theorem 2.1: Existence and Uniqueness of Mild Solutions}
These conditions come together in Theorem 2.1.

\begin{theorem}[Theorem 2.1]
Assume that the conditions \textnormal{(A.1)} through \textnormal{(A.4)} hold, together with the Lipschitz condition \textnormal{(H.1)} and the linear growth condition \textnormal{(H.2)}. Then there exists a unique mild solution to the SPDE. Moreover, this solution satisfies
\[
\mathbb P\left(\int_0^T \|X_s\|_H^2\,ds<\infty\right)=1.
\]
It also has a continuous modification, and for every \(p\ge 2\),
\[
\mathbb E\left(\sup_{t\in[0,T]}\|X_t\|_H^p\right)
\le
c_{p,T}\bigl(1+\mathbb E(\|\xi\|_H^p)\bigr).
\]
Thus the size of the solution over the whole time interval is controlled by the initial condition and the length of the interval.
\end{theorem}

Thus the mild solution is not only unique, but also has finite energy over the interval \([0,T]\) almost surely. This confirms that the infinite-dimensional stochastic construction is both robust and mathematically meaningful.

\section*{5. Proof Strategy for Theorem 2.1}
We now turn from the statement of Theorem 2.1 to the fixed point argument that proves it. Since the path regularity and maximal bound have already been established in the theorem, the remaining task is to construct the solution map and show that it is a contraction.

\subsection*{5.1. The Banach Fixed Point Framework}
The proof proceeds by Banach's Fixed Point Theorem. Instead of solving the SPDE directly, we build a Banach space of processes, define the mild solution operator on that space, and show that it becomes a contraction.

We begin by fixing a power \(p\ge 2\) and defining the function space \(H_p\):
\[
H_p
:=
\left\{
Y:\Omega_T\to H:
Y \text{ is predictable and }
\|Y\|_p:=\sup_{t\in[0,T]}\bigl(\mathbb E(\|Y(t)\|_H^p)\bigr)^{1/p}<\infty
\right\}.
\]
This space consists of all predictable, \(H\)-valued processes whose supremum of the expected \(p\)-th moment over time is finite. Equipped with this norm, \(H_p\) is a complete Banach space.

With this space fixed, for any process \(Y\in H_p\) we define the operator \(K\) by
\[
K(Y)(t)
:=
T_t(\xi)+\int_0^t T_{t-s}(B(s,Y(s)))\,ds+\int_0^t T_{t-s}(C(s,Y(s)))\,dW(s).
\]
To analyze it more smoothly, we decompose the operator into three parts:
\[
K(Y)(t)=T_t(\xi)+K_1(Y)(t)+K_2(Y)(t).
\]
Here \(K_1\) denotes the deterministic drift integral and \(K_2\) denotes the stochastic diffusion integral. The immediate goal is to prove that \(K\) is well defined, namely that it maps \(H_p\) back into \(H_p\).

\subsection*{5.2. Bounding the Drift and Diffusion Components}
With the operator \(K\) defined, we next show that it maps \(H_p\) into itself. We first estimate the drift component.
By applying H\"older's inequality to the Lebesgue time integral, we obtain a factor \(t^{p-1}\) and arrive at
\[
\mathbb E\left[\left\|\int_0^t T_{t-s}(B(s,Y(s)))\,ds\right\|_H^p\right]
\le
t^{p-1}\,
\mathbb E\left[\int_0^t \|T_{t-s}(B(s,Y(s)))\|_H^p\,ds\right].
\]
Using the uniform boundedness of the \(C_0\)-semigroup together with the linear growth condition \textnormal{(H.2)} (see Appendix A.16), this leads to the bound
\[
\mathbb E\bigl[\|K_1(Y)(t)\|_H^p\bigr]
\le
M(T)\,T^p\bigl(1+\|Y\|_p^p\bigr).
\]

We then turn to the diffusion component \(K_2\) and estimate its \(p\)-th moment.
Because this is an It\^o integral, we use the Burkholder-Davis-Gundy inequality:
\[
\mathbb E\left[\left\|\int_0^t T_{t-s}(C(s,Y(s)))\,dW(s)\right\|_H^p\right]
\le
c_p\,\mathbb E\left[\left(\int_0^t \|T_{t-s}(C(s,Y(s)))\circ \sqrt Q\|_{L_2(U,H)}^2\,ds\right)^{p/2}\right].
\]
Again, by the uniform boundedness of the semigroup and the linear growth condition on \(C\) (see Appendix A.17), we obtain
\[
\mathbb E\bigl[\|K_2(Y)(t)\|_H^p\bigr]
\le
M(T)\,T^{p/2}\bigl(1+\|Y\|_p^p\bigr).
\]
Since both \(K_1\) and \(K_2\) are bounded by expressions involving \(1+\|Y\|_p^p\), and the initial condition is also controlled, we conclude that \(K(Y)\in H_p\). Thus the operator maps bounded processes back to bounded processes.

\subsection*{5.3. The Contraction Estimate}
To apply Banach's theorem, it remains to show that \(K\) is a contraction. For two processes \(Y_1,Y_2\in H_p\), we therefore study
\[
\mathbb E\bigl[\|K(Y_1)(t)-K(Y_2)(t)\|_H^p\bigr].
\]
Using the standard inequality
\[
|a+b|^p\le 2^{p-1}(|a|^p+|b|^p),
\]
we split the total difference into the drift and diffusion parts:
\[
\mathbb E\bigl[\|K(Y_1)(t)-K(Y_2)(t)\|_H^p\bigr]
\le
2^{p-1}\Bigl(
\mathbb E\bigl[\|K_1(Y_1)(t)-K_1(Y_2)(t)\|_H^p\bigr]
+
\mathbb E\bigl[\|K_2(Y_1)(t)-K_2(Y_2)(t)\|_H^p\bigr]
\Bigr).
\]
Applying the global Lipschitz condition \textnormal{(H.1)} to the two terms above, and using the same time estimates as before, yields
\[
\|K(Y_1)-K(Y_2)\|_p^p
\le
M(T^{p/2}+T^p)\|Y_1-Y_2\|_p^p.
\]
Here \(M\) absorbs the Lipschitz constants, the uniform bounds of the \(C_0\)-semigroup, and the Burkholder-Davis-Gundy constants. The contraction estimate becomes effective as soon as
\[
M(T^{p/2}+T^p)<1
\qquad \text{for \(T\) small.}
\]
Since \(p\ge 2\), both \(T^{p/2}\) and \(T^p\) tend to zero as \(T\to 0\). Hence one can choose \([0,T]\) short enough so that \(K\) is a contraction on \(H_p\).

\section*{6. Completing the Fixed Point Argument}
With the contraction estimate established, the conclusion follows directly from Banach's theorem.

\subsection*{6.1. Banach's Fixed Point Theorem}
Consequently, there exists a unique process \(Y^*\in H_p\) such that
\[
K(Y^*)=Y^*.
\]
When we substitute this fixed point into the definition of \(K\), we obtain exactly the mild integral equation of the SPDE. Therefore this uniquely determined fixed point is the unique mild solution.

The contraction estimate is obtained first on a sufficiently short time interval. Since the smallness condition depends only on the interval length, the same argument can be repeated on consecutive intervals of the same size; together with the global linear growth condition \textnormal{(H.2)}, this extends the unique mild solution to any prescribed finite time horizon.

\section*{Appendix}
\subsection*{A.1. The Predictable \(\sigma\)-Algebra}
\begin{definition}[Predictable \(\sigma\)-algebra]
Let \((\mathcal F_t)_{t\in[0,T]}\) be a filtration. The predictable \(\sigma\)-algebra \(\mathcal P_T\) on \(\Omega\times[0,T]\) is defined by
\[
\mathcal P_T
:=
\sigma\bigl(\{]s,t]\times F_s:0\le s<t\le T,\ F_s\in\mathcal F_s\}\cup\{\{0\}\times F_0:F_0\in\mathcal F_0\}\bigr).
\]
\end{definition}

\subsection*{A.2. The Space \(L_2^0\)}
\begin{definition}[The space \(L_2^0\)]
The space \(L_2^0\) consists of all continuous linear operators \(T:U\to H\) such that
\[
T\circ \sqrt Q\in L_2(U,H).
\]
Equivalently,
\[
\sum_{k=1}^{\infty}\|T(\sqrt Q\,e_k)\|_H^2<\infty
\]
for one, hence every, complete orthonormal system \((e_k)\) of \(U\).
We need this space because the stochastic integral is defined only for integrands whose composition with \(\sqrt Q\) is Hilbert-Schmidt, so \(L_2^0\) is exactly the natural operator space in which the integrand must take values.
\end{definition}

\subsection*{A.3. Localization}
\begin{remark}
Localization means that one first defines the stochastic integral on a smaller square-integrable class of processes and then extends it to a larger class by stopping or truncating the integrands. In the present setting, the relevant condition is
\[
\mathbb P\left(\int_0^T \|\Phi(s)\circ \sqrt Q\|_{L_2(U,H)}^2\,ds<\infty\right)=1.
\]
This allows pathwise finiteness almost surely, even when the corresponding expectation is not imposed in advance.
\end{remark}

\subsection*{A.4. The Closure of \(\mathcal E\)}
\begin{remark}
The closure \(\overline{\mathcal E}\) is taken with respect to the norm induced by the Wiener-It\^o isometry. Concretely, elementary processes are completed in the \(L^2\)-sense after composition with \(\sqrt Q\), and equivalent processes are identified if they agree almost everywhere with respect to the corresponding norm.
\end{remark}

\subsection*{A.5. Burkholder-Davis-Gundy Inequalities}
\begin{remark}
The Burkholder-Davis-Gundy inequalities compare moments of the path supremum of a martingale with moments of its quadratic variation. In the present setting they give bounds of the form
\[
\mathbb E\left[\sup_{t\in[0,T]}\|M_t\|_H^{2p}\right]
\le
C_p\,\mathbb E\bigl(\langle M\rangle_T^p\bigr).
\]
\end{remark}

\subsection*{A.6. Quadratic Variation}
\begin{definition}[Quadratic variation]
For a continuous \(H\)-valued martingale \(M\), the quadratic variation \(\langle M\rangle_t\) is the unique continuous increasing adapted process starting at zero such that
\[
\|M_t\|_H^2-\langle M\rangle_t
\]
is a martingale.
\end{definition}

\subsection*{A.7. The Fractional Sobolev Space \(W^{\alpha,2}([0,T],H)\)}
\begin{definition}[Fractional Sobolev space]
For \(\alpha<1/2\), the space \(W^{\alpha,2}([0,T],H)\) consists of all \(H\)-valued functions \(M\in L^2([0,T],H)\) such that
\[
\int_0^T\int_0^T
\frac{\mathbb E\|M_s-M_t\|_H^2}{|s-t|^{1+2\alpha}}
\,ds\,dt<\infty.
\]
The double-integral term by itself is only a seminorm, which we denote by
\[
[M]_{W^{\alpha,2}([0,T],H)}^2
:=
\int_0^T\int_0^T
\frac{\mathbb E\|M_s-M_t\|_H^2}{|s-t|^{1+2\alpha}}
\,ds\,dt.
\]
In other words, the fractional increment term alone is only a seminorm. It is only a seminorm because it vanishes on functions that are constant in time.
The corresponding full norm is
\[
\|M\|_{W^{\alpha,2}([0,T],H)}^2
:=
\int_0^T \|M(t)\|_H^2\,dt
+
[M]_{W^{\alpha,2}([0,T],H)}^2.
\]
\end{definition}

\subsection*{A.8. Symmetry of the Double Integral}
\begin{remark}
In the regularity proof, one uses the symmetry of the kernel
\[
(s,t)\mapsto \frac{\mathbb E\|M_s-M_t\|_H^2}{|s-t|^{1+2\alpha}}.
\]
Indeed, interchanging \(s\) and \(t\) does not change the numerator or the denominator, so
\[
F(s,t):=\frac{\mathbb E\|M_s-M_t\|_H^2}{|s-t|^{1+2\alpha}}
\]
satisfies \(F(s,t)=F(t,s)\). Therefore
\[
\int_0^T\int_0^T F(s,t)\,ds\,dt
=
2\int_0^T\int_0^t F(s,t)\,ds\,dt.
\]
This is because the square \([0,T]\times[0,T]\) is the union of the two congruent triangles
\[
\{0\le s<t\le T\}
\qquad\text{and}\qquad
\{0\le t<s\le T\},
\]
and the symmetry \(F(s,t)=F(t,s)\) shows that the two contributions are equal.
Applied to the stochastic integral estimate, this gives
\begin{align*}
\mathbb E[X_T]
&=
\int_0^T\int_0^T
\frac{\mathbb E\left[\int_{s\wedge t}^{s\vee t}\|\Phi(r)\|_{L_2^0}^2\,dr\right]}
{|s-t|^{1+2\alpha}}
\,ds\,dt \\
&=
2\int_0^T\int_0^t
\frac{\mathbb E\left[\int_s^t\|\Phi(r)\|_{L_2^0}^2\,dr\right]}
{(t-s)^{1+2\alpha}}
\,ds\,dt.
\end{align*}
On the region \(0\le s<t\le T\), one has
\[
s\wedge t=s,
\qquad
s\vee t=t,
\]
so the inner integral
\[
\int_{s\wedge t}^{s\vee t}\|\Phi(r)\|_{L_2^0}^2\,dr
\]
reduces exactly to
\[
\int_s^t\|\Phi(r)\|_{L_2^0}^2\,dr.
\]
Thus the factor \(2\) comes from the symmetry of the square domain, while the change from \(s\wedge t\) and \(s\vee t\) to \(s\) and \(t\) comes from restricting the integral to the triangle \(0\le s<t\le T\).
\end{remark}

\subsection*{A.9. The Singular Kernel Bound}
\begin{lemma}
Let \(\alpha<1/2\). Then there exists a constant \(C(\alpha,T)\) such that
\[
\int_r^T\int_0^r \frac{1}{(t-s)^{1+2\alpha}}\,ds\,dt
\le C(\alpha,T)
\qquad \text{for all } r\in[0,T].
\]
\end{lemma}

\begin{proof}
For fixed \(r\in[0,T]\), let
\[
I(r)
:=
\int_r^T\int_0^r \frac{1}{(t-s)^{1+2\alpha}}\,ds\,dt.
\]
We first compute the inner integral. For fixed \(t\in[r,T]\), make the change of variables
\[
u=t-s.
\]
Then
\[
du=-ds,
\qquad
s=0 \Longrightarrow u=t,
\qquad
s=r \Longrightarrow u=t-r.
\]
Hence
\[
\int_0^r \frac{1}{(t-s)^{1+2\alpha}}\,ds
=
\int_t^{t-r} u^{-1-2\alpha}(-du)
=
\int_{t-r}^{t} u^{-1-2\alpha}\,du.
\]
If \(\alpha=0\), then
\[
\int_{t-r}^{t} u^{-1}\,du
=
\log t-\log(t-r)
\]
and the resulting bound is again finite on \([r,T]\). We therefore focus on the case \(\alpha\neq 0\). Then
\[
\int_{t-r}^{t} u^{-1-2\alpha}\,du
=
\left[-\frac{1}{2\alpha}u^{-2\alpha}\right]_{u=t-r}^{u=t}
=
\frac{(t-r)^{-2\alpha}-t^{-2\alpha}}{2\alpha}.
\]
Therefore
\[
I(r)
=
\frac{1}{2\alpha}\int_r^T \bigl((t-r)^{-2\alpha}-t^{-2\alpha}\bigr)\,dt.
\]
We now estimate the two time integrals separately. First,
\[
\int_r^T (t-r)^{-2\alpha}\,dt
=
\int_0^{T-r} v^{-2\alpha}\,dv
=
\frac{(T-r)^{1-2\alpha}}{1-2\alpha}
\le
\frac{T^{1-2\alpha}}{1-2\alpha},
\]
where we used the change of variables \(v=t-r\). Second,
\[
0\le \int_r^T t^{-2\alpha}\,dt \le \int_0^T t^{-2\alpha}\,dt<\infty,
\]
because \(-2\alpha>-1\) when \(\alpha<1/2\). Combining these estimates yields
\[
I(r)
\le
\frac{1}{2\alpha}\left(
\frac{T^{1-2\alpha}}{1-2\alpha}
+
\int_0^T t^{-2\alpha}\,dt
\right)
=:
C(\alpha,T)
\]
for all \(r\in[0,T]\). Thus the bound is uniform in \(r\), as claimed.
\end{proof}

\subsection*{A.10. Cylindrical Wiener Process}
\begin{remark}
A cylindrical Wiener process is formally written as
\[
W_t=\sum_{n\in\mathbb N}\beta_t^{(n)}u_n.
\]
Unlike a \(Q\)-Wiener process, this series does not converge in the original Hilbert space \(U\) in general. It should be understood as an idealized noise acting through admissible integrands.
\end{remark}

\subsection*{A.11. \(C_0\)-Semigroup}
\begin{definition}[\(C_0\)-semigroup]
A family \((T_t)_{t\ge 0}\subset L(H)\) is called a \(C_0\)-semigroup if
\[
T_0=\operatorname{Id},
\qquad
T_tT_s=T_{t+s},
\]
and for every \(u\in H\), the map \(t\mapsto T_tu\) is continuous.
\end{definition}

\subsection*{A.12. Infinitesimal Generator}
\begin{definition}[Generator]
If \((T_t)_{t\ge 0}\) is a \(C_0\)-semigroup on \(H\), its generator \(A\) is defined by
\[
Ah=\lim_{t\to 0}\frac{T_th-h}{t}
\]
on the domain
\[
D(A)=\left\{h\in H:\lim_{t\to 0}\frac{T_th-h}{t}\text{ exists in }H\right\}.
\]
\end{definition}

\subsection*{A.13. Basic Process Terminology}
\begin{remark}
Let \((X_t)_{t\in[0,T]}\) be a stochastic process.

\textbf{Uniquely determined process.} A process is uniquely determined if any two candidates satisfying the same defining properties agree in the relevant sense. In SPDE theory this is often meant up to indistinguishability or up to \(\mathbb P\)-almost sure equality for every \(t\).

\textbf{Continuous process.} The process is continuous if, for almost every \(\omega\), the map
\[
t\mapsto X_t(\omega)
\]
is continuous on \([0,T]\).

\textbf{Increasing process.} A real-valued process \((A_t)_{t\in[0,T]}\) is increasing if, for almost every \(\omega\),
\[
s\le t \quad \Longrightarrow \quad A_s(\omega)\le A_t(\omega).
\]
Thus each sample path is nondecreasing in time.

\textbf{\((\mathcal F_t)\)-adapted process.} The process is adapted to the filtration \((\mathcal F_t)\) if for every \(t\in[0,T]\), the random variable \(X_t\) is \(\mathcal F_t\)-measurable. This means that \(X_t\) depends only on the information available up to time \(t\).
\end{remark}

\subsection*{A.14. The Space \(H_p\)}
\begin{remark}
For \(p\ge 2\), the space \(H_p\) consists of predictable \(H\)-valued processes \(Y\) such that
\[
\sup_{t\in[0,T]}\bigl(\mathbb E(\|Y(t)\|_H^p)\bigr)^{1/p}<\infty.
\]
The norm appearing here is the natural \(p\)-moment supremum norm over time.
\end{remark}

\subsection*{A.15. Banach's Fixed Point Theorem}
\begin{theorem}[Banach's Fixed Point Theorem]
Let \((E,d)\) be a complete metric space and let \(K:E\to E\) be a contraction, that is,
\[
d(Kx,Ky)\le q\,d(x,y)
\]
for some \(q\in(0,1)\) and all \(x,y\in E\). Then \(K\) has a unique fixed point in \(E\).
\end{theorem}

\subsection*{A.16. The Drift Bound from Semigroup Boundedness and Linear Growth}
\begin{lemma}
Let \(p\ge 2\), let \((T_t)_{t\ge 0}\) be a uniformly bounded \(C_0\)-semigroup on \([0,T]\), and assume the linear growth condition \textnormal{(H.2)}. Then there exists a constant \(M(T)\) such that for every \(Y\in H_p\) and every \(t\in[0,T]\),
\[
\mathbb E\bigl[\|K_1(Y)(t)\|_H^p\bigr]
\le
M(T)\,T^p(1+\|Y\|_p^p).
\]
\end{lemma}
\begin{proof}
Assume that the \(C_0\)-semigroup is uniformly bounded on \([0,T]\), so that
\[
M_T:=\sup_{0\le r\le T}\|T_r\|_{L(H)}<\infty.
\]
We use the linear growth condition \textnormal{(H.2)} in the form
\[
\|B(t,x)\|_H^2+\|C(t,x)\circ \sqrt Q\|_{L_2(U,H)}^2\le M(1+\|x\|_H^2).
\]
Since \(p\ge 2\), this implies, after increasing the constant if necessary, that
\[
\|B(t,x)\|_H^p+\|C(t,x)\circ \sqrt Q\|_{L_2(U,H)}^p\le M_p(1+\|x\|_H^p).
\]
\noindent For the drift term, we start from
\[
\mathbb E\left[\left\|\int_0^t T_{t-s}(B(s,Y(s)))\,ds\right\|_H^p\right]
\le
t^{p-1}\,
\mathbb E\left[\int_0^t \|T_{t-s}(B(s,Y(s)))\|_H^p\,ds\right].
\]
Here we use H\"older's inequality for the time integral in the form
\[
\left(\int_0^t a(s)\,ds\right)^p
\le
t^{p-1}\int_0^t a(s)^p\,ds
\]
with \(a(s)=\|T_{t-s}(B(s,Y(s)))\|_H\).
By uniform boundedness of the semigroup,
\[
\|T_{t-s}(B(s,Y(s)))\|_H^p
\le
M_T^p\|B(s,Y(s))\|_H^p.
\]
Indeed, from
\[
M_T:=\sup_{0\le r\le T}\|T_r\|_{L(H)}<\infty
\]
we have \(\|T_{t-s}\|_{L(H)}\le M_T\) for every \(0\le s\le t\le T\). Hence
\[
\|T_{t-s}h\|_H\le \|T_{t-s}\|_{L(H)}\|h\|_H\le M_T\|h\|_H
\qquad\text{for all } h\in H,
\]
and applying this with \(h=B(s,Y(s))\) gives the displayed estimate.
Hence
\[
\mathbb E\left[\left\|\int_0^t T_{t-s}(B(s,Y(s)))\,ds\right\|_H^p\right]
\le
t^{p-1}M_T^p\int_0^t \mathbb E\bigl(\|B(s,Y(s))\|_H^p\bigr)\,ds.
\]
Using the growth bound on \(B\),
\[
\mathbb E\bigl(\|B(s,Y(s))\|_H^p\bigr)
\le
M_p\bigl(1+\mathbb E\|Y(s)\|_H^p\bigr)
\le
M_p(1+\|Y\|_p^p).
\]
Therefore
\[
\mathbb E\bigl[\|K_1(Y)(t)\|_H^p\bigr]
\le
t^pM_T^pM_p(1+\|Y\|_p^p)
\le
M(T)\,T^p(1+\|Y\|_p^p).
\]
\end{proof}

\subsection*{A.17. The Diffusion Bound from Semigroup Boundedness and Linear Growth}
\begin{lemma}
Let \(p\ge 2\), let \((T_t)_{t\ge 0}\) be a uniformly bounded \(C_0\)-semigroup on \([0,T]\), and assume the linear growth condition \textnormal{(H.2)}. Then there exists a constant \(M(T)\) such that for every \(Y\in H_p\) and every \(t\in[0,T]\),
\[
\mathbb E\bigl[\|K_2(Y)(t)\|_H^p\bigr]
\le
M(T)\,T^{p/2}(1+\|Y\|_p^p).
\]
\end{lemma}
\begin{proof}
Assume again that
\[
M_T:=\sup_{0\le r\le T}\|T_r\|_{L(H)}<\infty.
\]
We use the linear growth condition \textnormal{(H.2)} in the form
\[
\|B(t,x)\|_H^2+\|C(t,x)\circ \sqrt Q\|_{L_2(U,H)}^2\le M(1+\|x\|_H^2),
\]
which for \(p\ge 2\) implies
\[
\|C(t,x)\circ \sqrt Q\|_{L_2(U,H)}^p\le M_p(1+\|x\|_H^p).
\]
For the diffusion term, the Burkholder-Davis-Gundy inequality gives
\[
\mathbb E\left[\left\|\int_0^t T_{t-s}(C(s,Y(s)))\,dW(s)\right\|_H^p\right]
\le
c_p\,\mathbb E\left[\left(\int_0^t \|T_{t-s}(C(s,Y(s)))\circ \sqrt Q\|_{L_2(U,H)}^2\,ds\right)^{p/2}\right].
\]
Again, by uniform boundedness,
\[
\|T_{t-s}(C(s,Y(s)))\circ \sqrt Q\|_{L_2(U,H)}
\le
M_T\|C(s,Y(s))\circ \sqrt Q\|_{L_2(U,H)}.
\]
This uses the same bound \(\|T_{t-s}\|_{L(H)}\le M_T\), now at the Hilbert-Schmidt level: for every Hilbert-Schmidt operator \(S\in L_2(U,H)\),
\[
\|T_{t-s}\circ S\|_{L_2(U,H)}
\le
\|T_{t-s}\|_{L(H)}\|S\|_{L_2(U,H)}
\le
M_T\|S\|_{L_2(U,H)}.
\]
Applying this with \(S=C(s,Y(s))\circ \sqrt Q\) gives the displayed inequality.
So
\[
\mathbb E\left[\left\|\int_0^t T_{t-s}(C(s,Y(s)))\,dW(s)\right\|_H^p\right]
\le
c_pM_T^p\,
\mathbb E\left[\left(\int_0^t \|C(s,Y(s))\circ \sqrt Q\|_{L_2(U,H)}^2\,ds\right)^{p/2}\right].
\]
Since \(p/2\ge 1\), H\"older's inequality in time yields
\[
\left(\int_0^t a(s)^2\,ds\right)^{p/2}
\le
t^{\frac p2-1}\int_0^t a(s)^p\,ds.
\]
Applying this with \(a(s)=\|C(s,Y(s))\circ \sqrt Q\|_{L_2(U,H)}\), we obtain
\[
\mathbb E\left[\left(\int_0^t \|C(s,Y(s))\circ \sqrt Q\|_{L_2(U,H)}^2\,ds\right)^{p/2}\right]
\le
t^{\frac p2-1}\int_0^t \mathbb E\bigl(\|C(s,Y(s))\circ \sqrt Q\|_{L_2(U,H)}^p\bigr)\,ds.
\]
Using the growth bound on \(C\),
\[
\mathbb E\bigl(\|C(s,Y(s))\circ \sqrt Q\|_{L_2(U,H)}^p\bigr)
\le
M_p(1+\|Y\|_p^p).
\]
Therefore
\[
\mathbb E\bigl[\|K_2(Y)(t)\|_H^p\bigr]
\le
c_pM_T^pM_p\,t^{p/2}(1+\|Y\|_p^p)
\le
M(T)\,T^{p/2}(1+\|Y\|_p^p).
\]
\end{proof}

\end{document}
